\documentclass{article}
\usepackage[UTF8]{ctex}
\usepackage{amsmath,amssymb,amsfonts}
\usepackage{enumitem}
\usepackage{geometry}
\geometry{a4paper, margin=1in}
% === 错题本排版优化（不参与题目标识） ===
\setlength{\parindent}{0pt}
\setlength{\parskip}{0.35em plus 0.10em minus 0.08em}
\setlength{\emergencystretch}{3em}
\setlist{itemsep=0.25em, topsep=0.40em, parsep=0pt}
\allowdisplaybreaks[2]
% === 错题本排版优化结束 ===
\begin{document}
\section*{武汉理工大学2024年数学分析真题}
\section*{135．武汉理工大学2024年数学分析试卷}

\subsection{题目1}
1．计算极限 \(\lim _{n \rightarrow \infty}\left(e^{2}+6^{n}+8^{n}\right)^{\frac{1}{n}}\) ．

\subsection{题目2}
2．设 \(f(x)=\left(x^{2}+a\right) e^{x}\) 在定义域 \(\mathbb{R}\) 上有拐点但无极值点，求 \(a\) 的取值范围．

\subsection{题目3}
3．计算积分 \(\int_{0}^{\frac{\pi}{2}} \frac{\cos x}{\sin x+\cos x} \mathrm{~d} x\) ．

\subsection{题目4}
4．计算曲面积分

\[
\iint_{\sum} x^{3} \mathrm{~d} y \mathrm{~d} z+y^{3} \mathrm{~d} z \mathrm{~d} x+\left(x^{2}+y^{2}+z^{3}\right) \mathrm{d} x \mathrm{~d} y
\]

其中 \(\sum: x^{2}+y^{2}+z^{2}=R^{2}\) ，方向取外侧。

\subsection{题目5}
5．讨论级数 \(\sum_{n=1}^{\infty} \frac{x^{n}}{n(n+1)}\) 在 \([-1,1]\) 上的一致收敛性，并求其和函数．

\subsection{题目6}
6．设 \(f(x)\) 在 \(x=0\) 的某邻域内存在一阶连续导数，且 \(\lim _{x \rightarrow 0} \frac{f(x)}{x}=a>0\) 。讨论 \(\sum_{n=1}^{\infty}(-1)^{n} f\left(\frac{1}{n}\right)\) 是否收敛，若收敛是条件收敛还是绝对收敛？

\subsection{题目7}
7．已知 \(z=f(x, y)\) 的全微分 \(\mathrm{d} z=2 x \mathrm{~d} x-2 y \mathrm{~d} y\) ，且 \(f(1,1)=2\) ：

\begin{itemize}
  \item[（1）] 求 \(f(x, y)\) 的表达式．
  \item[（2）] 求 \(f(x, y)\) 在 \(D:\left\{(x, y) \left\lvert\, x^{2}+\frac{y^{2}}{4} \leq 1\right.\right\}\) 上的最大值与最小值．
\end{itemize}

\subsection{题目8}
8．讨论函数

\[
f(x, y)= \begin{cases}\frac{x y}{\sqrt{x^{2}+y^{2}}}, & (x, y) \neq(0,0) \\ 0, & (x, y)=(0,0)\end{cases}
\]

在原点处的连续性与可微性。

\subsection{题目9}
9．设 \(f(x)\) 存在一阶连续导数，且

\[
f(x)=x+\int_{0}^{x} t f^{\prime}(x-t) \mathrm{d} t+\int_{0}^{+\infty} e^{-\sqrt{x}} \mathrm{~d} x
\]

求 \(f(x)\) 的表达式．

\subsection{题目10}
10．证明如下问题：

\begin{itemize}
  \item[（1）] \(\frac{x}{1+x}<\ln (1+x)<x(x>0)\) ．
  \item[（2）] \(a_{n}=\sum_{k=1}^{n} \frac{1}{k}-\ln n\) 收敛．
\end{itemize}

\subsection{题目11}
11．用有限覆盖定理证明康拓定理，即 \([a, b]\) 上的连续函数一致连续．

\subsection{题目12}
12．设 \(f(x)\) 在 \([-a, a]\) 上二阶连续可导，且 \(f(0)=0\) ，证明：存在 \(\xi \in(-a, a)\) ，使得

\[
f^{\prime \prime}(\xi)=\frac{1}{a^{2}}[f(a)+f(-a)]
\]

\subsection{题目13}
13．已知连续函数列 \(\left\{f_{n}(x)\right\}\) 在 \([a, b]\) 上一致收敛于 \(f(x)\) ，若 \(\left\{x_{n}\right\} \subset[a, b]\) 满足 \(\lim _{n \rightarrow \infty} x_{n}=c\) ，证明：

\[
\lim _{n \rightarrow \infty} f_{n}\left(x_{n}\right)=f(c)
\]
\end{document}
